\section{Affine motion on neighborhoods of boxes}\label{ba:geometry}

Rectangle separation controls the beginning and end of a transfer; distinct
heights separate its middle. We now specify the paths and cutoff margins
used by the applications. Interpolating one planar scale and its reciprocal
preserves thickness and, for a linear first scale, gives a convex formula
for each first coordinate. Interpolating both planar scales independently
gives that formula in both coordinates, at the cost of a compensating
normal strain during the motion.

\subsection{Reciprocal paths and their localization}
\label{sec:ba-localization}

Use the flat step $\sigma$ of \eqref{eq:p2-step} and put
\[
 \theta(s)=\sigma(2s-1/2),\qquad
 w_j(s)=\theta(3s-j+1)\quad(j=1,2,3).
\]
The switches $w_1,w_2,w_3$ perform the lift, horizontal transfer and
descent in separate thirds of the unit interval. All are constant near
their phase endpoints. The following version of
Lemma~\ref{lem:bb-lift-parking}(i) records a common vertical half-width
of one and explicit neighborhoods; the same formulas allow other heights
and thicknesses as described immediately afterward.

\begin{lemma}[Two explicit interpolations of separated boxes]
\label{ba:two-interpolations}
Let $S_i,T_i\subset\mathbb R^2$, $1\le i\le N$, be closed rational
rectangles of positive side lengths.  Suppose each list is pairwise
disjoint and
\[
 F_i(y)=c_i^T+\operatorname{diag}(\lambda_i,\lambda_i^{-1})(y-c_i^S),
 \qquad\lambda_i\in\mathbb Q_{>0},
\]
maps $S_i$ onto $T_i$, with $c_i^S,c_i^T$ their centers.  There is a
smooth nondecreasing switch $\eta:[0,1]\to[0,1]$, zero near 0 and one
near 1, for which either scale
path $r_i=\lambda_i^\eta$ or $r_i=1-\eta+\eta\lambda_i$ yields an
effective smooth compactly supported solenoidal unit-time field, zero
on time collars, realizing $(y,z)\mapsto(F_i(y),z)$ on the whole box
$S_i\times[-1,1]$ and a neighborhood of it.  The paths are affine;
each horizontal half-width remains between its endpoint values, and the
horizontal sweep of each original box lies in the coordinate bounding
rectangle of $S_i$ and $T_i$. With the linear choice, a particle's first
coordinate is the convex
interpolation of its endpoint first coordinates.
\end{lemma}

\begin{proof}\label{ba:two-interpolations-proof}
For $N=0$ take the zero field.  For $N\ge1$ apply
Lemma~\ref{lem:bb-lift-parking}(i) with vertical half-width one.
Use the shared phase switches with $\eta=w_2$ and set
\[
 C_i(s)=\big((1-\eta)c_i^S+\eta c_i^T,\,6i(w_1-w_3)\big),
 \qquad D_i(s)=\operatorname{diag}(r_i,r_i^{-1},1).
\]
The affine path is
$G_i(s,y)=C_i(s)+D_i(s)(y-(c_i^S,0))$.
Choose $\epsilon$ to be one quarter of the minimum of 1 and all
pairwise max-norm distances in the two planar lists, omitting empty
pair lists.  The $\epsilon$-padded moving boxes are disjoint: source
and target gaps are at least $4\epsilon$ during the vertical phases,
and the middle-phase height spacing is $6>2(1+\epsilon)$.

For later loading formulas, record the centered cutoff explicitly:
\[
 \chi(C,L,\epsilon;x)=
 \prod_{j=1}^3\sigma\!\left(\frac{L_j+\epsilon+x_j-C_j}{\epsilon/2}\right)
 \sigma\!\left(\frac{L_j+\epsilon-x_j+C_j}{\epsilon/2}\right).
\]
Here $L$ is the vector of half-widths; the plateau has padding
$\epsilon/2$ and the support has padding $\epsilon$.
If $L_i(s)$ are the moving half-widths, $Y=x-C_i(s)$ and
$A_i=D_i'D_i^{-1}$, the field supplied by
\eqref{eq:bb-potential} is
\begin{equation}\label{ba:affine-box-field}
 V(s,x)=\sum_i\nabla\times\left\{
 \chi(C_i,L_i,\epsilon;x)
 \left[\tfrac12 C_i'\times Y+\tfrac13(A_iY)\times Y\right]
 \right\}.
\end{equation}
Indeed $\det D_i=1$, and the preceding separation verifies
Lemma~\ref{lem:bb-curl} with localization parameter $\epsilon/2$.
Its neighborhood conclusion applies with
$L_* =\max_i\max(1,\lambda_i,\lambda_i^{-1})$ because
$\|D_i(s)\|_\infty\le L_*$ and $D_i(0)=I$.
It therefore gives the stated affine motion on the full boxes and
on their initial max-norm neighborhoods of radius $\epsilon/(4L_*)$.
The paths are stationary for $s\le1/12$ and $s\ge11/12$.
Their finite center ranges and endpoint width bounds give common
compact support.  The effective positive lower bound
$r_i\ge\min(1,\lambda_i)$ verifies the scale requirement.
For the linear scale, the first-coordinate identity is
\[
 (G_i(s,y))_1=(1-\eta)y_1+\eta(F_i(y_1,y_2))_1.
\]
The second coordinate need not interpolate its own endpoint values.
Instead, if $b_{i2}$ is the source's second half-width, convexity of the
reciprocal function gives
\begin{equation}\label{eq:ba-reciprocal-envelope}
 \frac{b_{i2}}{(1-\eta)+\eta\lambda_i}
 \le (1-\eta)b_{i2}+\eta\frac{b_{i2}}{\lambda_i}.
\end{equation}
The first half-width is exactly its endpoint interpolation. For the
exponential scale, convexity of the exponential gives the corresponding
upper bounds in both coordinates, since
$\lambda_i^{\pm\eta}\le(1-\eta)+\eta\lambda_i^{\pm1}$.
In either case the center is interpolated linearly and each half-width
is at most the interpolation of its endpoint half-widths. Therefore each
lower edge is at least the interpolation of the endpoint lower edges,
and each upper edge is at most the corresponding interpolation of upper
edges. This proves the claimed bounding rectangle for the entire sweep.
The individual endpoint-width bounds also follow directly from monotonicity
of the positive scales and their reciprocals.
\end{proof}

The same construction applies to $P_i\times[-h,h]$ for any common
$h\ge0$, including plates and rectangles with zero side lengths, under
the hypotheses of Lemma~\ref{lem:bb-lift-parking}(i). Replace $6i$ by
heights $H_i$ whose pairwise gaps exceed $2h+2\varepsilon$, and choose
$\varepsilon$ below one third of every within-family planar gap. The
formula for $C_i$ uses the source and target centers in place of
$c_i^S,c_i^T$; its support padding is $\varepsilon$ and plateau padding
is $\varepsilon/2$. Source projections, distinct heights and target
projections separate the three phases in that order. Empty pair lists
impose no gap restriction, and the empty family gives the zero field.
Arbitrary real data give smooth existence; the effective-data assumptions
of the shared lemma give computable scales, derivative bounds and margins.
The same perturbation argument gives an open neighborhood even when a
side length is zero.

\subsection{Transporting an initial cutoff}
\label{sec:ba-transported-localization}

The transported-initial-cutoff convention is also useful.  Write the
affine path as $G_i(t)x=J_i(t)x+g_i(t)$ and localize by
$\psi_i(G_i(t)^{-1}x)$, where $\psi_i$ has support padding
$\varepsilon_0$ around the initial box and equals one on a smaller
positive padding.
If $\lambda_i\in[B^{-1},B]$, its initial padding
$\varepsilon_0$ expands horizontally by at most $B$.  Choose
$2B\varepsilon_0$ below every target gap,
$2\varepsilon_0$ below every source gap, and
$\varepsilon_0\le\varepsilon$.  These transported supports stay
separated in the same three phases.  The affine generator is
$v_i+A_ix$, with $A_i=\dot J_iJ_i^{-1}$ and
$v_i=\dot g_i-A_ig_i$.  Its potential is
\[
 \tfrac12v_i\times x+\tfrac13(A_ix)\times x.
\]
The identity proved in Lemma~\ref{lem:bb-curl} gives the exact
generator on the transported plateau.  Its neighborhood argument then
proves the same full-box endpoint map.  The factor $B$ in the margin
choice is essential: transported padding does not stay fixed under
stretching.

\subsection{Independent planar interpolation}

For an observer that needs both planar coordinates to follow their own
endpoint interpolations, the reciprocal path is insufficient. The next
construction interpolates the two factors independently and compensates
the temporary planar area change in the third direction. Here the named
targets may be disjoint containers of the actual affine images; only the
image centers enter the motion.

\begin{lemma}[Independent planar scales and normal compensation]
\label{ba:plate-routing}
Let $P_i,Q_i\subset\R^2$ be finite families of closed rectangles,
separately pairwise disjoint, and let assigned positive diagonal affine
maps $F_i$ satisfy $F_i(P_i)\subseteq Q_i$. Write their diagonal factors as
$\lambda_{i1}^*,\lambda_{i2}^*>0$ with product one.
There is a smooth compactly supported solenoidal unit-time field, zero on
time collars, whose endpoint map is $(y,z)\mapsto(F_i(y),z)$ on a
neighborhood of each entire plate $P_i\times\{0\}$. During the horizontal
phase both planar factors interpolate linearly, while the normal factor
compensates their product. Each planar coordinate of a point on the plate
is its convex endpoint interpolation.
\end{lemma}
\begin{proof}\label{ba:plate-routing-proof}
Use the height phases in Section~\ref{sec:ba-localization}. For source
center $p_i$, end its center path at its actual image $F_i(p_i)$; in the
onto case this is the named target center. Choose separated heights $H_i$
as above and set
\[
 c_i=\bigl((1-w_2)p_i+w_2F_i(p_i),\,H_i(w_1-w_3)\bigr),
\]
\[
 \lambda_{ij}=1-w_2+w_2\lambda_{ij}^*,\qquad
 D_i=\operatorname{diag}\bigl(
 \lambda_{i1},\lambda_{i2},(\lambda_{i1}\lambda_{i2})^{-1}\bigr).
\]
The volume-preserving path has logarithmic derivative
\[
 T_i=\dot D_iD_i^{-1}=\operatorname{diag}\left(
 \frac{\dot\lambda_{i1}}{\lambda_{i1}},
 \frac{\dot\lambda_{i2}}{\lambda_{i2}},
 -\frac{\dot\lambda_{i1}}{\lambda_{i1}}
 -\frac{\dot\lambda_{i2}}{\lambda_{i2}}\right).
\]
This matrix has trace zero, so \eqref{eq:bb-potential} localized
around the moving plate gives the desired field.  Its planar motion is
$c_i+(\lambda_{i1}(y_1-p_{i1}),\lambda_{i2}(y_2-p_{i2}),0)$.
Normal displacement changes by the factor
$(\lambda_{i1}\lambda_{i2})^{-1}$. The factors have positive
lower bounds; choose the initial thickness less than their product's
minimum times the normal plateau width.  A small enlargement of the
planar rectangle is allowed by the positive collars as well.  Thus a
whole thin neighborhood follows the volume-preserving affine motion.
The normal factor returns to one at the endpoint. The final plate image
is contained in $Q_i\times\{0\}$. Source and target-container
gaps protect the vertical phases and distinct heights protect the
middle phase.  Linear interpolation of each planar factor gives both
convex-coordinate identities on the plate.
This proves smooth existence for arbitrary real data. Under the effective
input assumptions of Lemma~\ref{lem:bb-curl}, the positive scale bounds and
finite gap margins can be computed and the construction is effective.
\end{proof}

\subsection{Forcing and effective evaluation}

\begin{lemma}[Residual forcing for the finite processors]\label{ba:realization}
Fix computable $\nu>0$.  Let $W$ be smooth, solenoidal, initially zero,
with support in one compact subset of $\mathbb R^3$ and every mixed
derivative bounded.  Then
\begin{equation}\label{ba:residual}
 F_W=W_t+(W\cdot\nabla)W-\nu\Delta W
\end{equation}
has the same support and derivative properties.  The zero-data solution
$(W,0)$ is unique among smooth solutions with
$u\in C_tH^2\cap C_t^1L^2$, bounded $u$ on finite intervals, and
$p\in C_tH^1$.  The corresponding torus assertion holds in the
classical comparison class of Section~\ref{sec:model}.  A sum of curls
in a torus chart gives both $W$ and $F_W$ mean zero.  Finite formulas
from the effective-data instances of the preceding constructions, finite
loading, and repetition of one cycle give effective evaluation of every
force derivative.
\end{lemma}
\begin{proof}\label{ba:realization-proof}
Fixed compact support and bounded mixed derivatives imply all reference
velocity hypotheses of Lemma~\ref{lem:b-comparison}.  Case (i) gives
exactly the stated whole-space comparison class, without a competing
gradient bound; its torus clause gives the stated classical comparison
class.  The same lemma provides the pressure normalization and global
material flow.  The product rule applied to \eqref{ba:residual}
proves the force's support and derivative bounds.  A chart curl has
zero mean, and integration of the residual uses
$\int\Delta W=0$ and
$\int(W\cdot\nabla)W=\int\operatorname{div}(W\otimes W)=0$.
For these effective-data instances, gate derivatives and seam evaluation are
those of \eqref{eq:p2-step} and Proposition~\ref{prop:bb-realization}.
The explicit affine scales and every reciprocal denominator have positive
computable bounds.  Thus differentiating the finite formulas, using a
finite superset of possible period indices, evaluates every derivative
without a machine execution or a real equality test at a seam.
\end{proof}
